\section{Appendix}

\begin{frame}{Appendix A: MLB Semantics}
\hypertarget{appendix:mlb-defs}{}
\begin{columns}[T,onlytextwidth]
\begin{column}{0.46\textwidth}
\begin{itemize}
    \item Let $\Sigma \subseteq G \times G \times G$ be a set of \ac{mlb} constraints.
    \item $MLB_{x y z}=(d_x,d_y,d_z)$ means $d_x$ and $d_y$ are closer to each other than to $d_z$.
\end{itemize}

\vspace{0.25em}
\textbf{Implication side}
\[
  d_x,d_z \in C \Rightarrow d_y \in C
\]

\textbf{Separation side}
\[
  \exists C:\ d_x,d_y \in C \land d_z \notin C
\]
\end{column}
\begin{column}{0.5\textwidth}
\banksearchHierarchyToMLB{}
\end{column}
\end{columns}
\end{frame}

\begin{frame}{Appendix B: Why $\phi_{\Sigma}$ Is a Closure Operator}
\hypertarget{appendix:closure-proof}{}
\begin{itemize}
    \item A set $A \subseteq G$ is $\Sigma$-consistent iff every constraint in $\Sigma$ is respected:
    \[
      \forall(d_x,d_y,d_z)\in\Sigma:\ \{d_x,d_z\}\subseteq A \Rightarrow d_y\in A .
    \]
    \item Let $\mathcal{S}_{\Sigma}$ be the family of all $\Sigma$-consistent sets.
    \item $G \in \mathcal{S}_{\Sigma}$, and arbitrary intersections of $\Sigma$-consistent sets are again $\Sigma$-consistent.
    \item Therefore $\mathcal{S}_{\Sigma}$ is a closure system.
\end{itemize}

\[
  \phi_{\Sigma}(A)
  =
  \bigcap \{ S \in \mathcal{S}_{\Sigma} \mid A \subseteq S \}
\]

\begin{block}{Consequence}
$\phi_{\Sigma}$ is extensive, monotone, and idempotent; hence it is a closure operator.
\end{block}
\end{frame}

\begin{frame}{Appendix C: Structural Saturation}
\hypertarget{appendix:saturation}{}
\begin{itemize}
    \item Standard \ac{fca} closure works on flat implications.
    \item A hierarchy can imply constraints that are not explicitly listed.
    \item Structural saturation closes this gap before $\phi_{\Sigma}$ is applied.
\end{itemize}

\vspace{0.35em}
\begin{center}
\begin{tikzpicture}[
    node distance=1.25cm,
    box/.style={draw,rounded corners=2pt,align=center,minimum width=2.7cm,minimum height=0.75cm},
    arr/.style={-{Latex[length=2mm]},thick}
]
\node[box] (i1) {$a,b \Rightarrow c$};
\node[box,below=0.45cm of i1] (i2) {$b,c \Rightarrow d$};
\node[box,right=1.8cm of $(i1)!0.5!(i2)$] (i3) {$a,b \Rightarrow d$};
\coordinate (i3upper) at ($(i3.west)!0.25!(i3.north west)$);
\coordinate (i3lower) at ($(i3.west)!0.25!(i3.south west)$);
\draw[arr] (i1.east) -- (i3upper);
\draw[arr] (i2.east) -- (i3lower);
\end{tikzpicture}
\end{center}

\vspace{0.35em}
\begin{block}{Interpretation}
Saturation makes indirect hierarchy information explicit, so the FCA closure operator can process it as ordinary implications.
\end{block}
\end{frame}


\begin{frame}{Step 4a: Ask for Exact Agreement}

\begin{itemize}
    \item Exact agreement means a BankSearch document cluster is closed under both knowledge sources.
    \item The corresponding closure operator is computed by alternating $\phi_{\Sigma}$ and $\phi_{TM}$ to a fixpoint.
    \item This produces the consistency lattice.
\end{itemize}
\end{frame}

\begin{frame}{BankSearch Agreement Matrix}
  \hypertarget{appendix:agreement-heatmap}{}
\begin{columns}[T,onlytextwidth]
\begin{column}{0.43\textwidth}
\begin{itemize}
    \item Most expert/data-driven concept pairs have near-zero overlap.
    \item The non-trivial agreements are concentrated in a few recognizable domains.
    \item This motivates a coherence lattice instead of stopping at consistency.
\end{itemize}
\end{column}
\begin{column}{0.55\textwidth}
\centering
\includegraphics[width=\textwidth,height=0.68\textheight,keepaspectratio]{heatmap_0.05_svg-tex.pdf}
\end{column}
\end{columns}
\end{frame}

\begin{frame}{Exact Agreement Is Rare}
\begin{columns}[T,onlytextwidth]
\begin{column}{0.50\textwidth}
\includegraphics[width=\textwidth]{heatmap_0.05_svg-tex.pdf}
\end{column}
\begin{column}{0.47\textwidth}
\begin{itemize}
    \item Similarity is computed between all concept extents.
    \item Apart from the top concept, most pairs have near-zero Jaccard similarity.
    \item Only five non-trivial pairs exceed $J(A,B)>0.5$.
    \item The manual hierarchy and data-driven explanation do not induce the same closed document sets.
\end{itemize}
\end{column}
\end{columns}
\end{frame}

\begin{frame}{Case Study as the Story Resolution TODO details}
\begin{itemize}
    \item Corpus: 11,000 documents.
    \item Represented topics: 11 categories organized into Finance, Science, Programming, and Sport.
    \item Expert hierarchy from prior BankSearch annotations is mapped to document-level \ac{mlb} constraints.
    \item Mapping yields 81 \ac{mlb} constraints.
\end{itemize}

\vspace{0.5em}
\begin{block}{Question revisited}
Can the manual hierarchy and the explainable topic-model \ac{fca} hierarchy support the same useful search clusters?
\end{block}
\end{frame}



\begin{frame}{Step 3: Build the Explainable Data-Driven Hierarchy TODO}
  \hypertarget{appendix:topic-model-details}{}
\begin{itemize}
    \item Train \ac{lda} with $T=11$ topics to match the represented BankSearch categories.
    \item Convert document-topic probabilities into a binary context with threshold $\delta_t=0.07$.
    \item Apply an iceberg support threshold $\delta_{ms}=0.05$ to retain interpretable frequent concepts.
    \item The resulting \ac{fca} concepts are explained by their topic intents.
\end{itemize}

\begin{center}
\includegraphics[height=0.37\textheight]{incidence_density_vs_threshold_svg-tex.pdf}
\hfill
\includegraphics[height=0.37\textheight]{concepts_vs_support_svg-tex.pdf}
\end{center}
\end{frame}